\documentclass[10pt]{article}
\input{../../notes-style.tex}

\usepackage{multicol}
\usepackage{array}

\title{Ch.~9 Notes\\\large AVL and Red-Black Trees}
\author{}
\date{}

\begin{document}
\maketitle
\vspace{-2em}

\begin{multicols}{2}

\subsection*{Why balanced BSTs}
Raw BST degrades to $O(n)$ on sorted input. AVL and RB are
BSTs that enforce a balance invariant so height stays
$O(\log n)$ \emph{worst case}.

\subsection*{AVL: the invariant}
For every node $n$, \textbf{balance factor}
\[
\mathrm{bf}(n) = h(n.\mathrm{left}) - h(n.\mathrm{right}) \in \{-1, 0, +1\}.
\]
Convention: height of empty subtree is $-1$; single node is $0$.
AVL worst-case height $\approx 1.44 \log_2(n+2)$.

\subsection*{AVL node}
\begin{lstlisting}[language=C++]
struct AVLNode {
    int key;
    int height = 0;     // cached
    AVLNode* left  = nullptr;
    AVLNode* right = nullptr;
    AVLNode* parent = nullptr;
};
\end{lstlisting}
Update \texttt{height} after every structural change:
$h = 1 + \max(h_L, h_R)$ using $-1$ for null children.

\subsection*{Rotations (BST-order-preserving swap)}
\begin{itemize}
    \item \textbf{Right-rotate at $y$}: $y$'s left child $x$ becomes
    root; $y$ becomes $x$'s right child; $x$'s old right subtree
    becomes $y$'s new left.
    \item \textbf{Left-rotate}: mirror.
\end{itemize}
In-order sequence is unchanged. This is the universal balancing
primitive.

\subsection*{The four AVL imbalance cases}
First unbalanced ancestor $z$; $y$ = taller child; $x$ = child of $y$
on the same side:
\begin{tabular}{@{}l|l@{}}
\textbf{shape} & \textbf{fix}\\
\hline
LL (left, left)  & right-rotate at $z$\\
RR (right, right)& left-rotate at $z$\\
LR (left, right) & left-rotate at $y$, then right-rotate at $z$\\
RL (right, left) & right-rotate at $y$, then left-rotate at $z$
\end{tabular}

\subsection*{AVL insert}
\begin{enumerate}
    \item BST insert at leaf; new node has height 0.
    \item \textbf{Retrace}: walk up from the new node, updating
    heights. At the \emph{first} ancestor with $|\mathrm{bf}| > 1$,
    apply the matching rotation. One rotation (single or double) is
    enough -- subtree height returns to its pre-insert value.
\end{enumerate}
Cost: $O(\log n)$ descent + $O(\log n)$ retrace = $O(\log n)$.

\columnbreak

\subsection*{AVL remove}
\begin{enumerate}
    \item BST remove (two-children $\to$ copy successor key, remove
    successor -- which has $\le 1$ child).
    \item Retrace from the removed node's parent. Unlike insert,
    rebalancing may \textbf{cascade}: fixing one level can leave the
    next one unbalanced.
\end{enumerate}

\subsection*{Red-black: the five rules}
\begin{enumerate}
    \item Every node is red or black.
    \item The root is black.
    \item A red node's children are black (no two reds in a row).
    \item Every null (NIL) leaf is black.
    \item Every root-to-NIL path has the same number of black nodes
    (the \emph{black-height}).
\end{enumerate}
Together $\Rightarrow$ longest path $\le 2 \times$ shortest $\Rightarrow$
height $\le 2 \log_2(n+1)$.

\subsection*{RB insertion fix-up, case summary}
New node inserted \textbf{red}. If parent is black, done. Else
(parent red -- rule 3 violated), by uncle:
\begin{itemize}
    \item \textbf{Uncle red} (case 3): recolor parent \& uncle black,
    grandparent red. Recurse at grandparent.
    \item \textbf{Uncle black, zig-zag} (case 4): rotate parent to
    straighten, then fall into case 5.
    \item \textbf{Uncle black, straight line} (case 5):
    recolor, rotate grandparent. Done.
\end{itemize}
Then: root $\to$ black. At most $O(\log n)$ recolorings and
$O(1)$ rotations.

\subsection*{RB removal (the hard one)}
BST-remove first. Only removing a \emph{black} node is tricky --
it breaks rule 5. Track a ``double-black'' token and fix via
six cases (red sibling, black sibling + red-nephew variants,
black sibling + black nephews $\to$ push deficit up).
Worst case $O(\log n)$ with $O(1)$ amortized rotations.

\subsection*{AVL vs.\ red-black}
\begin{tabular}{@{}l|l|l@{}}
         & \textbf{AVL}     & \textbf{red-black}\\
\hline
height   & $\le 1.44 \log n$ & $\le 2 \log n$\\
search   & faster (tighter) & slightly slower\\
insert   & up to 2 rotations& $\le 2$ rotations\\
remove   & may cascade      & $\le 3$ rotations\\
used by  & specialized DBs  & \texttt{std::map}, Linux kernel
\end{tabular}
Rule of thumb: AVL if reads dominate; RB if writes dominate.

\subsection*{Top gotchas}
\begin{itemize}
    \item Forgetting to update cached heights after a rotation.
    \item Reading balance factor with null children as height $0$
    instead of $-1$.
    \item Applying rotations at the wrong pivot -- read the cases
    as ``name from grandparent down.''
    \item Treating remove-retrace as ``one rotation fixes it'' --
    it can cascade all the way to the root.
    \item In RB: forgetting to reset root to black after a case-3
    recolor propagates it upward.
\end{itemize}

\end{multicols}
\end{document}
